% =====================================================================
%  PRÉAMBULE — Banque 1000 exercices · Maths Tle C (OnBuch)
%  Version DEUX COLONNES — dense, type « manuel de prépa »
%  Compatible XeTeX / tectonic. Mêmes noms de macros que l'original.
% =====================================================================

\usepackage[french]{babel}
\usepackage{microtype}

% --- Mise en page DENSE deux colonnes ---
\usepackage[a4paper,top=1.5cm,bottom=1.4cm,left=1.3cm,right=1.3cm,
            headheight=13pt,headsep=0.22cm,footskip=0.5cm]{geometry}
\usepackage{fancyhdr}
\usepackage{multicol}
\setlength{\columnsep}{16pt}
\setlength{\columnseprule}{0.4pt}
\renewcommand{\columnseprulecolor}{\color{onbhair}}
\setlength{\parindent}{0pt}
\setlength{\parskip}{1.2pt}
\linespread{0.98}
\setlength{\emergencystretch}{3.5em}   % évite les débordements de texte en colonnes étroites
\tolerance=3000 \hbadness=10000        % accepte des lignes lâches plutôt que de déborder
% Autorise la coupure des formules INLINE aux relations / opérateurs binaires
% (uniquement quand la ligne déborde) : indispensable en colonnes étroites.
\binoppenalty=80
\relpenalty=70

% --- Maths & graphiques ---
\usepackage{amsmath,amssymb,amsthm,mathrsfs}
\usepackage{bm}
\usepackage{esvect}
\usepackage{graphicx}
\usepackage{tikz}
\usetikzlibrary{calc,intersections,arrows.meta,decorations.markings,patterns,shapes.geometric,backgrounds}
\usepackage{pgfplots}
\pgfplotsset{compat=1.18}
\usepackage{xcolor}
\usepackage{tcolorbox}
\tcbuselibrary{skins,breakable}

% --- divers ---
\usepackage{enumitem}
\usepackage{multirow}
\usepackage{array}
\usepackage{booktabs}
\usepackage{titlesec}
\usepackage{titletoc}
\usepackage{etoolbox}
\usepackage{adjustbox}
\usepackage{pdfpages}
\usepackage[hidelinks]{hyperref}

% Auto-réduction des tableaux (tableaux de signes/variations) trop larges
\AtBeginDocument{%
  \BeforeBeginEnvironment{tabular}{\begin{adjustbox}{max width=\linewidth}}%
  \AfterEndEnvironment{tabular}{\end{adjustbox}}%
}

% Displays très resserrés (gros gain en deux colonnes) ----------------
% + auto-réduction à la largeur de COLONNE des displays \[...\] trop larges
% (ce qui tient garde sa taille ; seules les équations trop larges sont
%  réduites proportionnellement → plus aucun débordement dans la gouttière).
\newsavebox{\eqbox}
\newsavebox{\imbox}
% Maths INLINE longues : réduites à la largeur de colonne seulement si trop larges.
\newcommand{\fitmath}[1]{%
  \sbox{\imbox}{\ensuremath{#1}}%
  \ifdim\wd\imbox>\linewidth
    \par\nobreak\vspace{1pt}\noindent
    \makebox[\linewidth][c]{\resizebox{\linewidth}{!}{\ensuremath{\displaystyle #1}}}%
    \par\nobreak\vspace{1pt}\noindent
  \else \ensuremath{#1}\fi}   % branche « tient » : maths normale (reste sécable)
\newcommand{\outputeq}{%
  \par\addvspace{\abovedisplayskip}%
  \noindent\makebox[\linewidth][c]{%
    \resizebox{\ifdim\wd\eqbox>\linewidth \linewidth\else \wd\eqbox\fi}{!}{\usebox{\eqbox}}}%
  \par\addvspace{\belowdisplayskip}%
}
\AtBeginDocument{%
  \small                              % corps en \small : + de densité, équations + étroites
  \setlength{\abovedisplayskip}{2.5pt plus 1pt minus 1pt}%
  \setlength{\belowdisplayskip}{2.5pt plus 1pt minus 1pt}%
  \setlength{\abovedisplayshortskip}{1pt plus 1pt}%
  \setlength{\belowdisplayshortskip}{1.5pt plus 1pt}%
  \setlength{\jot}{1.5pt}%
  \def\[{\begin{lrbox}{\eqbox}$\displaystyle}%
  \def\]{$\end{lrbox}\outputeq}%
  % align / align* : redirigés vers aligned dans le display auto-réductible
  \renewenvironment{align}{\[\begin{aligned}}{\end{aligned}\]}%
  \renewenvironment{align*}{\[\begin{aligned}}{\end{aligned}\]}%
}

% =====================================================================
%  PALETTE ONBUCH (inchangée)
% =====================================================================
\definecolor{onbucorange}{HTML}{E8742C}
\definecolor{onbucorangedark}{HTML}{C8541A}
\definecolor{onbucdark}{HTML}{2C3E50}
\definecolor{onbuclight}{HTML}{FDF6EC}
\definecolor{onbucgreen}{HTML}{27AE60}
\definecolor{onbucgreendark}{HTML}{1E8449}
\definecolor{onbucblue}{HTML}{2E5C8A}
\definecolor{onbucbluelight}{HTML}{D6E4F0}
\definecolor{onbucgrey}{HTML}{7F8C8D}
\definecolor{onbucyellow}{HTML}{F39C12}
\definecolor{onbucpurple}{HTML}{8E44AD}
\definecolor{onbucpink}{HTML}{E91E63}
\definecolor{onbucteal}{HTML}{16A085}
\definecolor{onbucbg}{HTML}{FFF9F0}
\definecolor{pagebg}{HTML}{FFFEF8}
\definecolor{onbhair}{HTML}{E2D8CB}

\hypersetup{
  pdftitle={Banque d'exercices --- Mathématiques Tle C (Cameroun)},
  pdfauthor={OnBuch},
  pdfsubject={1000 exercices corrigés type Baccalauréat --- 23 chapitres}
}

% =====================================================================
%  EN-TÊTE / PIED
% =====================================================================
\pagestyle{fancy}
\fancyhf{}
\renewcommand{\headrulewidth}{0pt}
\renewcommand{\footrulewidth}{0pt}
\fancyhead[L]{\footnotesize\textcolor{onbucorange}{\rule[0.5pt]{5pt}{5pt}}\,\textcolor{onbucdark}{\bfseries Maths Tle C}}
\fancyhead[R]{\footnotesize\textcolor{onbucgrey}{\nouppercase{\leftmark}}}
\fancyfoot[C]{\footnotesize\textcolor{onbucgrey}{\thepage}}
\fancypagestyle{plain}{\fancyhf{}\renewcommand{\headrulewidth}{0pt}\fancyfoot[C]{\footnotesize\textcolor{onbucgrey}{\thepage}}}

% =====================================================================
%  TITRES
% =====================================================================
\titleformat{\chapter}[display]
  {\normalfont\bfseries}
  {\tikz[baseline=(n.base)]\node[circle,fill=onbucorange,text=white,inner sep=5pt,minimum size=1.05cm,font=\Large\bfseries] (n) {\thechapter};}
  {12pt}
  {\LARGE\color{onbucdark}}
  [\vspace{2pt}{\color{onbucorange}\titlerule[1.4pt]}]
\titlespacing*{\chapter}{0pt}{2pt}{10pt}

\titleformat{\section}
  {\normalfont\normalsize\bfseries\color{onbucorangedark}}
  {\colorbox{onbucorange}{\textcolor{white}{\,\thesection\,}}}{6pt}{}
\titlespacing*{\section}{0pt}{6pt}{3pt}
\titleformat{\subsection}
  {\normalfont\small\bfseries\color{onbucblue}}
  {\textcolor{onbucorange}{\thesubsection}}{5pt}{}
\titlespacing*{\subsection}{0pt}{5pt}{2pt}
% \paragraph en titre-BLOC compact (utilisé comme en-tête d'étape dans les corrigés)
\titleformat{\paragraph}[hang]
  {\normalfont\small\bfseries\color{onbucdark}}{}{0pt}{}
\titlespacing*{\paragraph}{0pt}{5pt}{2pt}

% =====================================================================
%  AUTO-WRAP DEUX COLONNES par chapitre (titre pleine largeur)
%  \chapter ferme les colonnes ouvertes, pose le titre pleine largeur,
%  puis rouvre deux colonnes. \chapter* (présentation) reste 1 colonne.
% =====================================================================
\newif\ifinmc \inmcfalse
\let\ORIchapter\chapter
\RenewDocumentCommand{\chapter}{s o m}{%
  \ifinmc\end{multicols}\inmcfalse\fi
  \IfBooleanTF{#1}
    {\ORIchapter*{#3}}%
    {\IfNoValueTF{#2}{\ORIchapter{#3}}{\ORIchapter[#2]{#3}}%
     \nobreak\begin{multicols}{2}\inmctrue}%
}
\AtEndDocument{\ifinmc\end{multicols}\inmcfalse\fi}

% =====================================================================
%  ÉNONCÉ / CORRIGÉ — en FLUX (circulent entre colonnes, jamais de débord)
% =====================================================================
% Énoncé : libellé non nécessaire (l'en-tête \titreexo le porte) ; léger fond ? non, flux pur.
\newenvironment{enoncebox}[1][]{\par\nobreak\ignorespaces}{\par}

% Corrigé : libellé vert en tête + filet vert fin sous le libellé.
\newenvironment{corrigeebox}[1][]{%
  \par\addvspace{1.5pt}\nobreak
  \noindent{\color{onbucgreendark}\bfseries\footnotesize\textsf{\textbullet\ CORRIGÉ}}%
  \par\nobreak\vspace{0.6pt}\noindent{\color{onbucgreen!55}\rule{\linewidth}{0.5pt}}\par\nobreak\vspace{1.2pt}%
  \ignorespaces}{\par\addvspace{1pt}}

% Contexte camerounais — court, garde une petite boîte teintée
\newtcolorbox{contextecmr}[1][]{%
  enhanced, breakable, frame hidden, colback=onbucbluelight!45,
  borderline west={1.4pt}{0pt}{onbucblue},
  left=6pt, right=3pt, top=1.2pt, bottom=1.2pt, boxsep=0pt,
  before skip=1.5pt, after skip=1.5pt, #1
}

% « L'essentiel du cours » — compact, en tête de chapitre (dans les colonnes)
\newtcolorbox{essentielbox}[1][]{%
  enhanced, breakable, colback=onbuclight!55, colframe=onbucpurple,
  boxrule=0pt, leftrule=2pt, arc=0.8mm,
  left=6pt, right=6pt, top=3pt, bottom=3pt,
  before skip=1pt, after skip=4pt,
  fonttitle=\bfseries\sffamily\footnotesize\color{white},
  coltitle=white, title={L'essentiel du cours},
  attach boxed title to top left={xshift=5pt,yshift=-2.2mm},
  boxed title style={colback=onbucpurple,arc=0.6mm,boxrule=0pt,size=fbox},
  #1
}

% Méthode — flux avec libellé teal
\newenvironment{methodebox}[1][]{%
  \par\addvspace{1.5pt}\nobreak
  \noindent{\color{onbucteal}\bfseries\footnotesize\textsf{\textbullet\ MÉTHODE}}%
  \par\nobreak\vspace{0.6pt}\noindent{\color{onbucteal!55}\rule{\linewidth}{0.5pt}}\par\nobreak\vspace{1.2pt}%
  \ignorespaces}{\par\addvspace{1pt}}

% =====================================================================
%  COMMANDES RACCOURCIES
% =====================================================================
\newcommand{\C}{\mathbb{C}}
\newcommand{\R}{\mathbb{R}}
\newcommand{\Z}{\mathbb{Z}}
\newcommand{\N}{\mathbb{N}}
\newcommand{\Q}{\mathbb{Q}}
\newcommand{\K}{\mathbb{K}}
\renewcommand{\i}{\mathrm{i}}
\DeclareMathOperator{\ch}{ch}
\DeclareMathOperator{\sh}{sh}
\renewcommand{\th}{\operatorname{th}}
\DeclareMathOperator{\argch}{argch}
\DeclareMathOperator{\argsh}{argsh}
\DeclareMathOperator{\argth}{argth}
\DeclareMathOperator{\pgcd}{pgcd}
\DeclareMathOperator{\ppcm}{ppcm}
\newcommand{\rep}[1]{\textcolor{onbucorangedark}{#1}}
\newcommand{\bareme}[1]{\hfill\textcolor{onbucgrey}{\footnotesize\textit{(#1\,pts)}}}

% Titre d'exercice COMPACT (numéro orange + intitulé + filet de séparation)
\newcommand{\titreexo}[2]{%
  \par\addvspace{5pt}\nobreak
  \noindent{\bfseries\color{onbucorangedark}Exercice~#1}%
  \hspace{5pt}{\footnotesize\color{onbucgrey}\textsc{#2}}%
  \par\nobreak\vspace{1pt}\noindent{\color{onbhair}\rule{\linewidth}{0.5pt}}%
  \par\nobreak\addvspace{1.5pt}%
}
\newcommand{\separateur}{\par\addvspace{3pt}}

% Théorèmes
\theoremstyle{definition}
\newtheorem{definition}{Définition}[section]
\newtheorem{propriete}{Propriété}[section]
\newtheorem{theoreme}{Théorème}[section]
\newtheorem{methodethm}{Méthode}[section]

% Listes compactes
\setlist[itemize]{itemsep=1pt,topsep=1pt,parsep=0pt,leftmargin=1.3em}
\setlist[enumerate]{itemsep=1pt,topsep=1pt,parsep=0pt,leftmargin=1.6em}
